\section{Simulation design}\label{sec:design}

The simulation uses a generator with known $\theta$ and varies one factor at a time around a base scenario, so that each theoretical result has a scenario that tests it against a prior numerical prediction.

\paragraph{Generator.}
\begin{enumerate}
  \item \emph{Covariates:} $V=(V_1,\dots,V_4)$ normal with correlations $0.3^{|j-k|}$, plus a three-level categorical $V_5$ whose distribution depends on $V_1$; all are customer characteristics. An unobserved $U$, independent of $V$, enters only the unobserved-confounding scenario.
  \item \emph{Latent channel variable:} $H=\ind\{C^*\ne\vacio\}$ indicates an actual history, with $\PP(H=1\mid V)=\sigma(-0.4+0.8V_1+0.3V_3)$; if $H=1$, the history takes the value $N(0.5V_1,1)$. It is recorded if $D=1$ and $H=1$, and $C=\vacio$ otherwise.
  \item \emph{Channel:} $\PP(D=1\mid V,H,U)=\sigma(\alpha_0+\lambda\,\alpha^\top V+\zeta H+\nu U)$. The parameter $\lambda$ controls overlap, and $\zeta\ne0$ makes the distribution of $C^*$ given $V$ differ between channels, a necessary condition for $C^*$ to confound.
  \item \emph{Outcomes:} two common products, $\PP(Y_k(d)=1\mid X^*,U)=\sigma(\beta_{k0}+\beta_{kd}+\beta_k^\top V+\delta H+\nu'U+e)$ for $k=1,2$, with $e$ a customer random effect, and a product exclusive to the premium channel with $Y_3(0)=0$. The outcome used by the estimators is $Y=Y_1+Y_2$.
  \item \emph{Target group:} $S=\ind\{D=1,\ s^*(V)\le c^*\}$, with $s^*$ the true purchase-propensity index in the premium channel and $c^*$ its 10th percentile in that channel. The rule is fixed, so $\theta$ does not change across replications.
  \item \emph{Units and observations:} each customer generates between 1 and 5 calls, with channel and covariates constant across them.
  \item \emph{Truth:} $\theta$ and the theoretical limit of each estimator are computed on a population of 1.5 million customers; the true conditional means and $dP_V/dQ_V$ are obtained by integrating out $H$, $U$ and the customer effect with Gauss--Hermite quadrature. Their Monte Carlo error, about 0.001 points (0.015 in the unobserved-confounding scenario), is negligible relative to that of the replication means. Trimming effects, much smaller, are predicted as paired differences within the population.
\end{enumerate}

\paragraph{Scenarios and predictions.} The base scenario excludes $C$, has $\zeta\ne0$ and $\delta=\nu=\nu'=0$, high overlap ($\lambda=0.5$; $\chi^2=3.6$), LightGBM with five-fold cross-fitting by customer for $\hat p$ and $\hat m$, and 30{,}000 customers per replication: about 62{,}000 donor calls and 2{,}800 target-group calls. Around it we vary the inclusion of $C$, confounding through the history ($\delta=1.2$), overlap ($\lambda=1.5$ and $3$; $\chi^2=6.3$ and 36), model specification, unobserved confounding ($\nu=\nu'=1$), trimming at the 99th, 95th and 90th percentiles of the donor weights, and size (3{,}000 to 30{,}000 customers). Each scenario has 200 replications, so that 95\,\% coverage is estimated with a standard error of about 1.5 points. Predictions that do not follow directly from the theory were registered, with a date, before running the corresponding scenario; the code, the predictions and the raw results are available.

\paragraph{The AUC does not measure overlap.} When calibrating $\lambda$, the AUC of the ideal classifier stayed between 0.90 and 0.92 for $\lambda$ between 0.3 and 2, while $\chi^2$ ranged from 3.6 to about 12. The definition of the target group, a cutoff on a score, already separates it from the donor channel. We therefore set the levels by $\chi^2$, which is what governs the variance (Proposition~\ref{prop:variance}).
